% Figure: mesh-convergence study for the peak stress at a hole. Peak stress
% (FEM) versus inverse element size, approaching the analytical value with
% Richardson extrapolation.
\begin{figure}[htbp]
\centering
\begin{tikzpicture}
\begin{axis}[
  width=11cm, height=6.5cm,
  xlabel={mesh refinement $1/h$ at the hole (1/mm)},
  ylabel={peak stress at hole (MPa)},
  xmin=0, xmax=11, ymin=120, ymax=160,
  axis lines=left,
  grid=major, grid style={enggray!20},
  tick label style={font=\scriptsize},
  label style={font=\small},
  legend style={font=\scriptsize}, legend pos=south east,
]
% FEM measured points: coarse under-predicts, refining approaches ~150
\addplot[engblue, very thick, mark=*, mark options={fill=engblue}] coordinates {
  (1, 128)
  (2, 139)
  (4, 146)
  (8, 149)
};
\addlegendentry{FEM peak stress}
% analytical / converged value
\addplot[engred, dashed, domain=0:11, samples=2] {150};
\addlegendentry{analytical $K_t\sigma = 150$ MPa}
% Richardson-extrapolated converged value marker
\addplot[englime, only marks, mark=square*, mark options={fill=englime}]
  coordinates {(10.5, 150)};
\node[englime, font=\scriptsize, anchor=east] at (axis cs:10.2,153)
  {Richardson limit};
\end{axis}
\end{tikzpicture}
\caption[Mesh convergence of peak stress]{A mesh convergence study
for the plate with a central hole project. The finite element peak
stress at the hole (blue) rises toward the analytical value of
$150\,\text{MPa}$ ($K_t = 3$ times the $50\,\text{MPa}$ remote
stress) as the mesh is refined. The coarse mesh under predicts the
peak by about $15\,\%$, which would have certified an unsafe design;
only after the study confirms convergence is the number an
engineering result. Richardson extrapolation (green square) recovers
the mesh independent value from the sequence of runs.}
\label{fig:vol03ch09:mesh-convergence}
\end{figure}
